\chapter{Introduction}
\label{ch:introduction}

\section{Motivation}
Finding an animal in a video and following it over time is a need shared by several fields. Ecologists use it
to survey wild populations from camera-trap footage. Neuroscientists and behavioural biologists use it to
quantify what an animal does frame by frame in the laboratory, where the tracked trajectory is often the
measurement the science rests on, and where a whole tool-chain now exists to recover pose and turn it into
behavioural labels~\cite{deeplabcut,sleap,poser}. The bottleneck is the same in both cases: video is cheap to
record and expensive to annotate.

This dissertation works on the wild side of that divide, where the problem is hardest. Camera traps have
become a standard tool in wildlife monitoring, and a single survey leaves a team with far more footage than
anyone can watch. Somebody still has to work out which animal is in each clip, how many there are, and what
they do. Done by hand this is slow and expensive, and it is the step that limits how much ecology gets done at
all.

Automating it means solving two problems, not one. The first is \emph{finding} the animal in a frame. The
second is \emph{following} it from frame to frame, so the same individual keeps the same identity for as long
as it is visible. Only the second yields the counts and behavioural sequences both communities want: in a
study of social behaviour, losing an identity means losing the very animal whose behaviour was being measured.
It is also the harder of the two.

Neither is easy on real footage. Animals are frequently camouflaged against the very background they live in.
They appear at any distance and from any viewpoint, so one species may fill the frame in one clip and occupy a
handful of pixels in the next. Vegetation cuts them in half. Much footage is recorded at night in infrared,
where colour is gone and contrast is poor. Motion blur is common. A box drawn around the animal is often a
poor description of it, which is why segmentation --- labelling the pixels that belong to the animal --- is
the more useful output. And because a camera trap is fixed, the same cluttered scene repeats in every clip
from that site. That gives a model an easy shortcut: learn the background rather than the
animal~\cite{terra,panaf,xiao}.

For years the standard answer was to train a detector on a fixed list of species and deploy it. It works where
it was trained. It degrades sharply at camera sites and regions it has not seen~\cite{terra,wilds,panaf}. And
adding a species that was not on the list means collecting and labelling more data --- precisely the cost a
team opening a new site is trying to avoid.

Promptable models change that arrangement. The \emph{Segment Anything} line replaced the fixed label set with a
prompt: you indicate what you want and the model segments it, and SAM\,2 extended this to
video~\cite{sam2}. SAM\,3~\cite{sam3} takes the step that matters here, because the prompt can now simply be a
phrase. Give it \texttt{"impala"} and it finds, segments and follows every impala in the clip, having never
been trained on that species. It is the current state of the art for the task, and it was released alongside
SA-FARI, a large camera-trap tracking benchmark of wild animals~\cite{safari}. A team can now point a tracker
at a species it has never seen and get an answer straight away.

That is the appeal. The liability is that the answer is not equally good everywhere. Zero-shot performance
varies sharply from one animal to the next, and from one place to the next. No rule tells a practitioner,
before running the model, whether its output can be trusted for a given animal in a given place. The choice
left is an unattractive one: trust the output blindly, or fall back on the manual annotation the model was
meant to replace.

This dissertation studies that gap. It asks whether such a rule exists --- first for SAM\,3, and then whether
whatever we find is specific to it.

\section{Problem statement and research question}
This dissertation asks whether a promptable tracker's zero-shot accuracy on an unseen (species, place) can be
predicted in advance from properties measurable without running the model or scoring its output, and which
factors govern
that transfer --- separately for finding the animal and following it. We study this concretely on
SAM\,3 as the primary case, measuring novelty against the SA-FARI reference split (Section~\ref{sec:bg-safari}),
and then test whether the answer generalises beyond one dataset and one model --- to two further datasets and,
via a controlled model swap, to two other promptable trackers (Section~\ref{sec:model_swap}).

\section{Hypotheses}
\begin{description}[leftmargin=2.6em,style=nextline]
  \item[H0 (null)] the label-free distances carry no predictive information about
    \texttt{pDetA}/\texttt{pAssA} --- transfer is not distance-driven.
  \item[H1 (detection)] \texttt{pDetA} falls with species novelty (taxonomic + visual distance).
  \item[H2 (association)] \texttt{pAssA} falls with environment difficulty (scene distance,
    day/night/IR, clutter).
\end{description}
Failing to reject H0 is itself an informative result --- and, as Chapter~\ref{ch:results} shows, it is
what we find: neither a novelty distance nor a follow-up intrinsic-difficulty pivot predicts
transfer beyond a trivial animal-size effect. We even read familiarity from SAM\,3's own features, and it too is
null; a fuller representational probe --- of what those features structurally encode --- remains the open
branch for future work.

\section{Contributions}
To our knowledge this is the first study to test, out of sample, whether the zero-shot transfer of a
promptable video tracker can be predicted before the model is run and from label-free distances --- and the
first to split that question into detection (\texttt{pDetA}) and association (\texttt{pAssA}). This distinguishes it
from label-free performance prediction for classifiers, which estimates a single scalar from the model's own
outputs on the target~\cite{aotl,agl,atc}, and from the very recent static-image detector/segmenter
estimators~\cite{pcr,rca}; and from camera-trap shift studies, which measure but do not forecast per-site
reliability~\cite{terra,panaf,jeon}. Concretely we deliver:
\begin{enumerate}[label=(\roman*),leftmargin=2.4em]
  \item a rigorous, out-of-sample negative result --- under a deliberately powerful test (whole-species
    and whole-location hold-out with group-cluster-bootstrap significance), label-free before-running distances
    do \emph{not} predict SAM\,3's transfer;
  \item the detection-versus-association decomposition, tested on two near-orthogonal splits --- a constructed
    species hold-out (novelty) and the official location hold-out (environment);
  \item a hardened account of why it fails: a follow-up intrinsic-difficulty pivot collapses to the
    same animal-size confound, so the only label-free property that predicts detection is object size --- the
    very nuisance the analysis controls for --- and that positive control confirms the test detects signal when
    it exists;
  \item evidence that the null is not an artefact of one dataset or one model: it holds across
    two datasets (SA-FARI and BURST) and, in a controlled model swap on BURST, across two
    architecturally distinct promptable trackers (GLEE and Florence-2\,+\,SAM\,2), so it is task-general rather
    than a quirk of SAM\,3.
\end{enumerate}

\section{Dissertation outline}
Chapter~\ref{ch:background} reviews the related work; Chapter~\ref{ch:methodology} describes the data,
measurement, distance features, and modelling; Chapter~\ref{ch:results} reports the results and
out-of-sample validation; Chapter~\ref{ch:discussion} discusses them; Chapter~\ref{ch:conclusion} concludes.
